\subsection*{组合：\texttt{outer \# inner}}

\begin{lstlisting}[language=croqtile,caption={组合：outer \# inner}]
output.at(tr # i, tc # j)
\end{lstlisting}

中缀 \co@#@ 将外层分块坐标和块内坐标组合成原张量中的坐标。

对于外层分块索引 \co@tr@ 和块内索引 \co@i@，\co@tr # i@ 的值为 \co@tr * #i + i@。每块有 16 行时，第 2 块的块内第 5 行对应全局第 37 行。组合后的上界是两个上界的乘积。内层上界决定步幅，因此操作数顺序不能交换。

图~\ref{fig:f04-fig_compose} 展示了分块坐标与块内坐标如何合成为全局坐标。

\begin{figure}[htbp]\centering
\includegraphics[width=\linewidth,height=0.56\textheight,keepaspectratio]{f04-fig_compose_light.png}
\caption{每块含 16 个元素时，tile = 2 与块内偏移 i = 3 组合得到全局位置 2 × 16 + 3 = 35。}
\label{fig:f04-fig_compose}
\end{figure}
